%% kapa: exam=2 round=32 year=2021 period=1 subject=appraisal_practice subject_ko=감정평가실무 seq=1 number=1 points=40
%% source: Q-Net 감정평가사 기출문제 (cst00309) artlSeq=5208993 file=raw/1008/5208993/감평 2차 시험지(2021).pdf
%% conversion: kordoc 4.24.51 -> tools/gen2.py (auto); review=proofread-sonnet; flags=none
\begin{longquestion}{【문제 1】}
감정평가사 甲은 S시에 소재하는 대상 부동산에 대하여 일반거래(시가 참고) 목적의 감정평가를 의뢰받았다. 관련법규 및 이론을 참작하고 제시된 자료를 활용하여 다음 각 물음에 답하시오. (40점) \par
\medskip\noindent 물음 1) 토지는 공시지가기준법, 거래사례비교법을 적용하고, 건물은 원가법을 적용하여 대상 부동산의 시산가액을 산정하시오. (18점) \par
\medskip\noindent 물음 2) 일괄 거래사례비교법에 의한 시산가액을 산정하시오. (7점) \par
\medskip\noindent 물음 3) 일괄 수익환원법에 의한 시산가액을 산정하시오. (12점) \par
\medskip\noindent 물음 4) 시산가액 조정을 통하여 감정평가액을 결정하시오. (3점) \par
\medskip\noindent < 자료 1 > 기본적 사항 \par
\medskip\noindent 1. 기준가치: 시장가치 \par
2. 기준시점: 2021년 8월 7일 \par
3. 대상 물건의 개황 \par
1) 토지 \par
{\small\centering
\begin{tabular}{|c|c|c|c|c|c|c|c|}\hline
\makecell{소재지\\지번} & 지목 & \makecell{면적\\(m\textsuperscript{2})} & \makecell{용도\\지역} & \makecell{이용\\상황} & \makecell{도로\\접면} & \makecell{형상\\지세} & 주위환경 \\\hline
\makecell{J구 M동\\120} & 대 & 1,500 & 일반상업 & 업무용 & \makecell{광대\\세각} & \makecell{가장형\\평지} & \makecell{일반\\업무지대} \\\hline
\end{tabular}\par}
2) 건물 \par
(1) 건물 개황 \par
{\small\centering
\begin{tabular}{|c|c|c|c|c|c|c|}\hline
\makecell{소재지\\지번} & 구조 & 층수 & \makecell{면적\\(m\textsuperscript{2})} & 용도 & 급수 & 비고 \\\hline
\makecell{J구 M동\\120} & \makecell{철근콘\\크리트조} & \makecell{지하4층/\\지상10층} & 13,800 & 업무용 & 3 & \makecell{허가일: 2015. 07. 15.\\사용승인일: 2016. 07. 15.\\(지상9\textasciitilde{}10층 증축:\\2018. 07. 15.)} \\\hline
\end{tabular}\par}
(2) 건물 세부 내역 \par
{\small\centering
\begin{tabular}{|c|c|c|c|}\hline
구분 & 면적(m\textsuperscript{2}) & 이용상황 & 부대설비 내역 \\\hline
지하1층\textasciitilde{}지하4층 & 각 950 & 주차장, 기계실 & \makecell{전기설비, 소방설비,\\승강기설비} \\\hline
지상1층\textasciitilde{}지상10층 & 각 1,000 & 업무시설 & \makecell{전기설비, 소방설비, 위생설비,\\냉난방설비, 승강기설비} \\\hline
\end{tabular}\par}
\medskip\noindent < 자료 2 > 공시지가표준지 \hfill (공시기준일: 2021.01.01) \par
{\small\centering
\begin{tabular}{|c|c|c|c|c|c|c|c|c|c|}\hline
\makecell{기\\호} & \makecell{소재지\\지번} & \makecell{지\\목} & \makecell{면적\\(m\textsuperscript{2})} & \makecell{용도\\지역} & \makecell{이용\\상황} & \makecell{도로\\접면} & \makecell{형상\\지세} & 주위환경 & \makecell{공시지가\\(원/m\textsuperscript{2})} \\\hline
1 & \makecell{J구 M동\\60} & 대 & 450 & \makecell{3종\\일주} & \makecell{상업\\용} & \makecell{중로\\한면} & \makecell{세장형\\평지} & \makecell{후면\\상가지대} & 22,000,000 \\\hline
2 & \makecell{J구 M동\\110} & 대 & 1,400 & \makecell{일반\\상업} & \makecell{업무\\용} & \makecell{광대\\한면} & \makecell{세장형\\평지} & \makecell{일반\\업무지대} & 41,000,000 \\\hline
3 & \makecell{J구 M동\\210} & 대 & 1,050 & \makecell{일반\\상업} & \makecell{업무\\용} & \makecell{소로\\한면} & \makecell{가장형\\평지} & \makecell{후면\\상가지대} & 30,000,000 \\\hline
\end{tabular}\par}
\medskip\noindent < 자료 3 > 인근지역 평가사례 및 거래사례 \par
1. 평가사례 \par
{\small\centering
\begin{tabular}{|c|c|c|c|c|c|c|c|c|c|c|}\hline
\makecell{기\\호} & \makecell{소재지\\지번} & \makecell{지\\목} & \makecell{면적\\(m\textsuperscript{2})} & \makecell{용도\\지역} & \makecell{이용\\상황} & \makecell{도로\\접면} & \makecell{형상\\지세} & \makecell{기준\\시점} & \makecell{토지단가\\(원/m\textsuperscript{2})} & \makecell{평가\\목적} \\\hline
가 & \makecell{J구 M동\\75} & 대 & 570 & \makecell{3종\\일주} & \makecell{상업\\용} & \makecell{중로\\한면} & \makecell{정방형\\평지} & \makecell{2021.\\02. 01.} & 38,500,000 & 담보 \\\hline
나 & \makecell{J구 M동\\105} & 대 & 1,300 & \makecell{일반\\상업} & \makecell{업무\\용} & \makecell{광대\\한면} & \makecell{세장형\\평지} & \makecell{2021.\\03. 01.} & 62,000,000 & \makecell{시가\\참고} \\\hline
다 & \makecell{J구 M동\\115} & 대 & 1,200 & \makecell{일반\\상업} & \makecell{업무\\용} & \makecell{광대\\한면} & \makecell{가장형\\평지} & \makecell{2018.\\06. 01.} & 58,000,000 & \makecell{시가\\참고} \\\hline
라 & \makecell{J구 M동\\125} & 대 & 1,400 & \makecell{일반\\상업} & \makecell{업무\\용} & \makecell{광대\\한면} & \makecell{가장형\\평지} & \makecell{2021.\\07. 01.} & 61,000,000 & \makecell{자산\\재평가} \\\hline
마 & \makecell{J구 M동\\195} & 대 & 1,360 & \makecell{일반\\상업} & \makecell{업무\\용} & \makecell{소로\\한면} & \makecell{가장형\\평지} & \makecell{2021.\\01. 01.} & 42,000,000 & 담보 \\\hline
\end{tabular}\par}
- 평가사례 기호 가, 마는 후면 상가지대, 기호 나 \textasciitilde{} 라는 일반 업무지대에 위치함 \par
\medskip\noindent 2. 거래사례 \par
(1) 거래사례 \#1 \par
- 소재지: J구 M동 109 \par
- 총 거래가격: 67,050,000,000원 \par
- 거래시점: 2021년 3월 1일 \par
- 토지: 일반상업, 주상용, 900m\textsuperscript{2}, 광대한면, 세장형, 평지 \par
- 건물 \par
{\small\centering
\begin{tabular}{|c|c|c|c|c|}\hline
구조 & 급수 & \makecell{연면적\\(m\textsuperscript{2})} & \makecell{허가일 /\\사용승인일} & 부대설비 내역 \\\hline
철근콘크리트조 & 4 & 12,500 & \makecell{2017. 02. 23. /\\2018. 02. 20.} & \makecell{전기설비, 소방설비, 위생설비,\\냉난방설비, 승강기설비} \\\hline
\end{tabular}\par}
- 기타사항: 일반 업무지대에 위치하며, 정상 거래사례임 \par
(2) 거래사례 \#2 \par
- 소재지: J구 M동 129 \par
- 총 거래가격: 98,400,000,000원 \par
- 거래시점: 2021년 2월 1일 \par
- 토지: 일반상업, 업무용, 1,600m\textsuperscript{2}, 광대세각, 세장형, 평지 \par
- 건물 \par
{\small\centering
\begin{tabular}{|c|c|c|c|c|}\hline
구조 & 급수 & \makecell{연면적\\(m\textsuperscript{2})} & \makecell{허가일 /\\사용승인일} & 부대설비 내역 \\\hline
철근콘크리트조 & 3 & 5,000 & \makecell{1980. 01. 20. /\\1981. 01. 25.} & \makecell{전기설비, 소방설비, 위생설비,\\냉난방설비} \\\hline
\end{tabular}\par}
- 기타사항: 일반 업무지대에 위치하는 정상적인 거래사례로, 매수자는 대상 부동산을 매입하여 지하4층, 지상10층 규모의 업무시설을 신축할 예정임(철거비는 감안하지 않는 것으로 함) \par
(3) 거래사례 \#3 \par
- 소재지: J구 M동 139 \par
- 총 거래가격: 99,636,000,000원 \par
- 거래시점: 2021년 3월 1일 \par
- 토지: 일반상업, 업무용, 1,500m\textsuperscript{2}, 광대한면, 가장형, 평지 \par
- 건물 \par
{\small\centering
\begin{tabular}{|c|c|c|c|c|}\hline
구조 & 급수 & \makecell{연면적\\(m\textsuperscript{2})} & \makecell{허가일 /\\사용승인일} & 부대설비 내역 \\\hline
철근콘크리트조 & 3 & 13,600 & \makecell{2015. 02. 16. /\\2016. 02. 19.} & \makecell{전기설비, 소방설비, 위생설비,\\냉난방설비, 승강기설비} \\\hline
\end{tabular}\par}
- 기타사항: 일반 업무지대에 위치하며, 매도자의 사정으로 인해 급매된 사례임 \par
(4) 거래사례 \#4 \par
- 소재지: J구 M동 209 \par
- 총 거래가격: 81,940,000,000원 \par
- 거래시점: 2021년 4월 1일 \par
- 토지: 일반상업, 업무용, 1,470m\textsuperscript{2}, 소로한면, 가장형, 평지 \par
- 건물 \par
{\small\centering
\begin{tabular}{|c|c|c|c|c|}\hline
구조 & 급수 & \makecell{연면적\\(m\textsuperscript{2})} & \makecell{허가일 /\\사용승인일} & 부대설비 내역 \\\hline
철근콘크리트조 & 3 & 11,000 & \makecell{2018. 03. 17. /\\2019. 03. 29.} & \makecell{전기설비, 소방설비, 위생설비,\\냉난방설비, 승강기설비} \\\hline
\end{tabular}\par}
- 기타사항: 후면 상가지대에 위치하며, K사의 펀드운용을 위한 투자목적으로 거래된 정상 거래사례임 \par
(5) 거래사례 \#5 \par
- 소재지: J구 M동 153 \par
- 총 거래가격: 111,573,000,000원 \par
- 거래시점: 2020년 10월 1일 \par
- 토지: 일반상업, 업무용, 1,600m\textsuperscript{2}, 광대한면, 가장형, 평지 \par
- 건물 \par
{\small\centering
\begin{tabular}{|c|c|c|c|c|}\hline
구조 & 급수 & \makecell{연면적\\(m\textsuperscript{2})} & \makecell{허가일 /\\사용승인일} & 부대설비 내역 \\\hline
철근콘크리트조 & 3 & 14,700 & \makecell{2015. 08. 20. /\\2016. 09. 20.} & \makecell{전기설비, 소방설비, 위생설비,\\냉난방설비, 승강기설비} \\\hline
\end{tabular}\par}
- 기타사항: 일반 업무지대에 위치하는 정상 거래사례임 \par
(6) 거래사례 \#6 \par
- 소재지: J구 M동 163 \par
- 총 거래가격: 102,900,000,000원 \par
- 거래시점: 2020년 9월 1일 \par
- 토지: 일반상업, 업무용, 1,500m\textsuperscript{2}, 광대한면, 가장형, 평지 \par
- 건물 \par
{\small\centering
\begin{tabular}{|c|c|c|c|c|}\hline
구조 & 급수 & \makecell{연면적\\(m\textsuperscript{2})} & \makecell{허가일 /\\사용승인일} & 부대설비 내역 \\\hline
철근콘크리트조 & 3 & 14,000 & \makecell{2014. 07. 16. /\\2015. 08. 19.} & \makecell{전기설비, 소방설비, 위생설비,\\냉난방설비, 승강기설비} \\\hline
\end{tabular}\par}
- 기타사항: 일반 업무지대에 위치하며, 대상 부동산은 구분소유건물로서 매수 후 개별분양 예정임 \par
(7) 거래사례 \#7 \par
- 소재지: J구 M동 173 \par
- 총 거래가격: 62,300,000,000원 \par
- 거래시점: 2020년 11월 1일 \par
- 토지: 일반상업, 업무용, 1,800m\textsuperscript{2}, 광대한면, 가장형, 평지 \par
- 건물 \par
{\small\centering
\begin{tabular}{|c|c|c|c|c|}\hline
구조 & 급수 & \makecell{연면적\\(m\textsuperscript{2})} & \makecell{허가일 /\\사용승인일} & 부대설비 내역 \\\hline
철근콘크리트조 & 4 & 10,000 & \makecell{2001. 09. 22. /\\2002. 11. 01.} & \makecell{전기설비, 소방설비, 위생설비,\\냉난방설비, 승강기설비} \\\hline
\end{tabular}\par}
- 기타사항: 일반 업무지대에 위치하며, 인근 중개업소에 탐문조사한 결과 거래에 따른 양도소득세는 매수자가 부담하는 것으로 조사되었음 \par
\medskip\noindent < 자료 5 > 재조달원가 및 감가수정 관련 자료 \par
1. 표준단가 \par
{\small\centering
\begin{tabular}{|c|c|c|c|c|}\hline
용도 & 구조 & 급수 & 표준단가(원/m\textsuperscript{2}) & 내용연수 \\\hline
업무시설 & \makecell{철근콘크리트조\\(6층\textasciitilde{}15층 이하)} & 1 & 1,400,000 & 50 \\\hline
업무시설 & \makecell{철근콘크리트조\\(6층\textasciitilde{}15층 이하)} & 2 & 1,300,000 & 50 \\\hline
업무시설 & \makecell{철근콘크리트조\\(6층\textasciitilde{}15층 이하)} & 3 & 1,200,000 & 50 \\\hline
업무시설 & \makecell{철근콘크리트조\\(6층\textasciitilde{}15층 이하)} & 4 & 1,100,000 & 50 \\\hline
업무시설 & \makecell{철근콘크리트조\\(6층\textasciitilde{}15층 이하)} & 5 & 1,000,000 & 50 \\\hline
\end{tabular}\par}
- 지상·지하 구분 없이 적용 가능함 \par
2. 부대설비 보정단가 \par
{\small\centering
\begin{tabular}{|c|c|}\hline
구분 & 보정단가(원/m\textsuperscript{2}) \\\hline
전기설비 & 10,000 \\\hline
소방설비 & 10,000 \\\hline
위생설비 & 50,000 \\\hline
냉난방설비 & 140,000 \\\hline
승강기설비 & 30,000 \\\hline
\end{tabular}\par}
3. 건물 잔가율은 0\%임 \par
4. 건물의 감가수정은 정액법(만년감가)을 적용함 \par
\medskip\noindent < 자료 6 > 시점수정 자료 \par
1. 지가변동률(S시 J구) \par
{\small\centering
\begin{tabular}{|c|c|c|}\hline
구분 & 주거지역 & 상업지역 \\\hline
2018. 06. 01.\textasciitilde{}2021. 06. 30.(누계) & 12.825 & 12.846 \\\hline
2020. 09. 01.\textasciitilde{}2021. 06. 30.(누계) & 4.057 & 4.036 \\\hline
2020. 10. 01.\textasciitilde{}2021. 06. 30.(누계) & 3.715 & 3.694 \\\hline
2020. 11. 01.\textasciitilde{}2021. 06. 30.(누계) & 3.376 & 3.355 \\\hline
2020. 12. 01.\textasciitilde{}2021. 06. 30.(누계) & 3.018 & 2.997 \\\hline
2021. 01. 01.\textasciitilde{}2021. 06. 30.(누계) & 2.624 & 2.645 \\\hline
2021. 02. 01.\textasciitilde{}2021. 06. 30.(누계) & 2.265 & 2.285 \\\hline
2021. 03. 01.\textasciitilde{}2021. 06. 30.(누계) & 1.827 & 1.845 \\\hline
2021. 04. 01.\textasciitilde{}2021. 06. 30.(누계) & 1.278 & 1.293 \\\hline
2021. 05. 01.\textasciitilde{}2021. 06. 30.(누계) & 0.795 & 0.806 \\\hline
2021. 06. 01.\textasciitilde{}2021. 06. 30. & 0.414 & 0.420 \\\hline
\end{tabular}\par}
- 2021년 7월 이후 지가변동률은 미고시 되었음 \par
2. 오피스빌딩 자본수익률(S시 J구) \par
{\small\centering
\begin{tabular}{|c|c|c|c|c|}\hline
구분 & 2020. 3분기 & 2020. 4분기 & 2021. 1분기 & 2021. 2분기 \\\hline
자본수익률(\%) & 0.42 & 0.46 & 0.50 & 0.54 \\\hline
\end{tabular}\par}
3. 건축비지수는 동일하다고 가정함 \par
\medskip\noindent < 자료 7 > 지역요인 \par
1. 대상과 공시지가표준지 및 사례는 인근지역에 소재하여 지역요인은 유사함 \par
\medskip\noindent < 자료 8 > 토지 개별요인 \par
1. 가로조건(각지인 경우 가로조건에서 반영하기로 함) \par
{\small\centering
\begin{tabular}{|c|c|c|c|c|c|c|c|}\hline
구분 & 광대한면 & 광대소각 & 광대세각 & 중로한면 & 중로각지 & 소로한면 & 소로각지 \\\hline
광대한면 & 1.00 & 1.09 & 1.05 & 0.95 & 0.99 & 0.85 & 0.89 \\\hline
광대소각 & 0.92 & 1.00 & 0.96 & 0.87 & 0.91 & 0.78 & 0.82 \\\hline
광대세각 & 0.95 & 1.04 & 1.00 & 0.90 & 0.94 & 0.81 & 0.85 \\\hline
중로한면 & 1.05 & 1.15 & 1.11 & 1.00 & 1.04 & 0.89 & 0.94 \\\hline
중로각지 & 1.01 & 1.10 & 1.06 & 0.96 & 1.00 & 0.86 & 0.90 \\\hline
소로한면 & 1.18 & 1.28 & 1.24 & 1.12 & 1.16 & 1.00 & 1.05 \\\hline
소로각지 & 1.12 & 1.22 & 1.18 & 1.07 & 1.11 & 0.96 & 1.00 \\\hline
\end{tabular}\par}
2. 접근조건 \par
{\small\centering
\begin{tabular}{|c|c|c|c|c|c|c|c|}\hline
구분 & 대상 & 표준지 & \makecell{평가사례\\가} & \makecell{평가사례\\나} & \makecell{평가사례\\다} & \makecell{평가사례\\라} & \makecell{평가사례\\마} \\\hline
평점 & 95 & 100 & 93 & 100 & 93 & 100 & 95 \\\hline
\end{tabular}\par}
\par{\small\centering
\begin{tabular}{|c|c|c|c|c|c|c|c|}\hline
구분 & \makecell{거래사례\\\#1} & \makecell{거래사례\\\#2} & \makecell{거래사례\\\#3} & \makecell{거래사례\\\#4} & \makecell{거래사례\\\#5} & \makecell{거래사례\\\#6} & \makecell{거래사례\\\#7} \\\hline
평점 & 93 & 95 & 93 & 100 & 90 & 95 & 93 \\\hline
\end{tabular}\par}
- 상기의 접근조건 비교치 산정시 소수점 셋째자리에서 반올림하여 소수점 둘째자리까지 산정함 \par
3. 획지조건 \par
{\small\centering
\begin{tabular}{|c|c|c|c|c|c|}\hline
구분 & 정방형 & 가장형 & 세장형 & 사다리형 & 부정형 \\\hline
정방형 & 1.00 & 1.00 & 0.98 & 0.95 & 0.92 \\\hline
가장형 & 1.00 & 1.00 & 0.98 & 0.95 & 0.92 \\\hline
세장형 & 1.02 & 1.02 & 1.00 & 0.97 & 0.94 \\\hline
사다리형 & 1.05 & 1.05 & 1.03 & 1.00 & 0.97 \\\hline
부정형 & 1.09 & 1.09 & 1.07 & 1.03 & 1.00 \\\hline
\end{tabular}\par}
4. 제시된 조건 외의 조건은 동일함 \par
\medskip\noindent < 자료 9 > 토지, 건물 일괄 개별요인 \par
1. 개별요인 \par
(1) 대상물건/거래사례 \#1 \par
{\small\centering
\begin{tabular}{|c|c|c|c|}\hline
구분 & 입지적 특성 & 기능적 특성 & 물리적 특성 \\\hline
대상물건 & 102 & 103 & 102 \\\hline
거래사례 \#1 & 100 & 100 & 100 \\\hline
\end{tabular}\par}
(2) 대상물건/거래사례 \#2 \par
{\small\centering
\begin{tabular}{|c|c|c|c|}\hline
구분 & 입지적 특성 & 기능적 특성 & 물리적 특성 \\\hline
대상물건 & 100 & 105 & 105 \\\hline
거래사례 \#2 & 100 & 100 & 100 \\\hline
\end{tabular}\par}
(3) 대상물건/거래사례 \#3 \par
{\small\centering
\begin{tabular}{|c|c|c|c|}\hline
구분 & 입지적 특성 & 기능적 특성 & 물리적 특성 \\\hline
대상물건 & 102 & 103 & 102 \\\hline
거래사례 \#3 & 100 & 100 & 100 \\\hline
\end{tabular}\par}
(4) 대상물건/거래사례 \#4 \par
{\small\centering
\begin{tabular}{|c|c|c|c|}\hline
구분 & 입지적 특성 & 기능적 특성 & 물리적 특성 \\\hline
대상물건 & 95 & 103 & 102 \\\hline
거래사례 \#4 & 100 & 100 & 100 \\\hline
\end{tabular}\par}
(5) 대상물건/거래사례 \#5 \par
{\small\centering
\begin{tabular}{|c|c|c|c|}\hline
구분 & 입지적 특성 & 기능적 특성 & 물리적 특성 \\\hline
대상물건 & 105 & 102 & 100 \\\hline
거래사례 \#5 & 100 & 100 & 100 \\\hline
\end{tabular}\par}
(6) 대상물건/거래사례 \#6 \par
{\small\centering
\begin{tabular}{|c|c|c|c|}\hline
구분 & 입지적 특성 & 기능적 특성 & 물리적 특성 \\\hline
대상물건 & 100 & 102 & 100 \\\hline
거래사례 \#6 & 100 & 100 & 100 \\\hline
\end{tabular}\par}
(7) 대상물건/거래사례 \#7 \par
{\small\centering
\begin{tabular}{|c|c|c|c|}\hline
구분 & 입지적 특성 & 기능적 특성 & 물리적 특성 \\\hline
대상물건 & 102 & 103 & 102 \\\hline
거래사례 \#7 & 100 & 100 & 100 \\\hline
\end{tabular}\par}
2. 상기의 개별요인은 상승식으로 계산하며, 제시된 특성 외의 특성은 동일함 \par
\medskip\noindent < 자료 10 > 대상부동산 및 인근지역 임대 현황 \par
1. 대상부동산 임대 현황 \par
{\small\centering
\begin{tabular}{|c|c|c|c|c|}\hline
구분 & \makecell{임대면적\\(m\textsuperscript{2})} & \makecell{월임대료\\(원/m\textsuperscript{2})} & \makecell{보증금\\(원/m\textsuperscript{2})} & \makecell{월관리비\\(원/m\textsuperscript{2})} \\\hline
지상1층 & 1,000 & 47,000 & 470,000 & 12,000 \\\hline
지상2층 & 1,000 & \multicolumn{3}{c|}{공실} \\\hline
지상3층\textasciitilde{}지상5층 & 3,000 & 25,000 & 250,000 & 12,000 \\\hline
지상6층\textasciitilde{}지상7층 & 2,000 & 35,000 & 350,000 & 12,000 \\\hline
지상8층\textasciitilde{}지상9층 & 2,000 & 27,000 & 270,000 & 12,000 \\\hline
지상10층 & 1,000 & 35,000 & 350,000 & 12,000 \\\hline
\end{tabular}\par}
- 지상3\textasciitilde{}5층과 지상8\textasciitilde{}9층은 각각 특수관계회사가 저가로 임차하고 있음 \par
2. 인근지역의 표준적 임대 현황(최근 자료) \par
{\small\centering
\begin{tabular}{|c|c|c|c|}\hline
구분 & 월임대료(원/m\textsuperscript{2}) & 보증금(원/m\textsuperscript{2}) & 월관리비(원/m\textsuperscript{2}) \\\hline
지상1층 & 47,000 & 470,000 & 14,000 \\\hline
지상2층\textasciitilde{}지상10층(각) & 35,000 & 350,000 & 14,000 \\\hline
\end{tabular}\par}
\medskip\noindent < 자료 11 > 수익환원법 적용 자료 \par
1. 인근지역 시장조사 결과 월임대료는 보증금의 10\% 수준으로 조사됨 \par
2. 보증금 운용이율은 연 2\%임 \par
3. 인근지역 건물의 전형적인 공실률은 5\%임 \par
4. 인근지역의 전형적인 운영경비는 관리비수입의 65\%임 \par
5. 인근지역에서 매년 1개월의 렌트프리(Rent Free)가 계약조건에 포함되는 것이 일반적인 시장관행임 \par
6. 인근지역 업무시설의 임대료는 전형적인 수준이 형성되어 있으며, 공실률 감소를 위한 유인책으로 임대계약시 랜트프리(Rent Free)를 적극적으로 활용하고 있음 \par
\medskip\noindent < 자료 12 > 환원이율 관련 자료 \par
1. 인근지역 유사부동산 자료(최근) \par
(1) 사례 \#101 \par
- 총 거래가격: 151,120,000,000원 \par
- 토지면적(용도지역): 1,960m\textsuperscript{2}(일반상업지역) \par
- 임대면적: 13,000m\textsuperscript{2} \par
- 임대현황 \par
{\small\centering
\begin{tabular}{|c|c|c|c|}\hline
보증금 & 관리비 & \makecell{랜트프리\\(Rent Free)} & \makecell{가능조소득\\(Potential Gross Income)} \\\hline
340,000원/m\textsuperscript{2} & 13,000원/m\textsuperscript{2} & 1개월 & 7,420,400,000원 \\\hline
\end{tabular}\par}
- 기타사항: 경매낙찰사례로, 최초 법사가격은 180,000,000,000원이며, 유치권 행사중임 \par
(2) 사례 \#102 \par
- 총 거래가격: 125,346,000,000원 \par
- 토지면적(용도지역): 1,887m\textsuperscript{2}(일반상업지역) \par
- 임대면적: 11,000m\textsuperscript{2} \par
- 임대현황 \par
{\small\centering
\begin{tabular}{|c|c|c|c|}\hline
보증금 & 관리비 & \makecell{랜트프리\\(Rent Free)} & \makecell{가능조소득\\(Potential Gross Income)} \\\hline
320,000원/m\textsuperscript{2} & 13,000원/m\textsuperscript{2} & 1개월 & 6,010,400,000원 \\\hline
\end{tabular}\par}
- 기타사항: 장기임차인의 임대재계약으로 저가 임대 중 거래된 사례임 \par
(3) 사례 \#103 \par
- 총 거래가격: 132,960,000,000원 \par
- 토지면적(용도지역): 1,695m\textsuperscript{2}(일반상업지역) \par
- 임대면적: 12,000m\textsuperscript{2} \par
- 임대현황 \par
{\small\centering
\begin{tabular}{|c|c|c|c|}\hline
보증금 & 관리비 & \makecell{랜트프리\\(Rent Free)} & \makecell{가능조소득\\(Potential Gross Income)} \\\hline
350,000원/m\textsuperscript{2} & 14,000원/m\textsuperscript{2} & 1개월 & 7,140,000,000원 \\\hline
\end{tabular}\par}
- 기타사항: 2인이 공유지분으로 소유권이전등기 된 사례임 \par
(4) 사례 \#104 \par
- 총 거래가격: 102,250,000,000원 \par
- 토지면적(용도지역): 2,070m\textsuperscript{2}(일반상업지역, 제3종일반주거지역) \par
- 임대면적: 10,000m\textsuperscript{2} \par
- 임대현황 \par
{\small\centering
\begin{tabular}{|c|c|c|c|}\hline
보증금 & 관리비 & \makecell{랜트프리\\(Rent Free)} & \makecell{가능조소득\\(Potential Gross Income)} \\\hline
300,000원/m\textsuperscript{2} & 14,000원/m\textsuperscript{2} & 1개월 & 5,340,000,000원 \\\hline
\end{tabular}\par}
- 기타사항: 정상 거래된 사례임 \par
2. 환원이율은 백분율로 소수점 둘째자리에서 반올림하여 백분율로 소수점 첫째자리까지 표시함 \par
3. 상기 각 사례의 보증금은 층별 임대면적을 가중평균하여 산정한 금액임 \par
\medskip\noindent < 자료 13 > 기타사항 \par
1. 공시지가표준지 및 사례 선정시 선정 및 제외 사유를 반드시 기재할 것 \par
2. < 자료 6 >을 이용한 시점수정치는 소수점 여섯째자리에서 반올림하여 소수점 다섯째자리까지 표시함 \par
3. 개별요인은 조건간 상승식으로 산정하되, 소수점 넷째자리에서 반올림하여 소수점 셋째자리까지 표시함 \par
4. 그 밖의 요인 보정치를 산정하는 경우 비교표준지를 기준으로 하는 방식을 적용함 \par
5. 공시지가기준법 및 거래사례비교법에 의한 토지 단가와 일괄 거래사례비교법에 의한 건물 면적당 단가는 반올림하여 각각 유효숫자 셋째자리까지 표시하며, 각 평가방법별 시산가액은 천만원 단위에서 반올림하여 억원 단위까지 표시함 \par
6. 일괄 거래사례비교법 적용시 건물 연면적을 기준으로 함 \par
\end{longquestion}
